%%%%%%%%%%%%%%%%%%%%%%%%%%%%%%%%%%%%%%%%%
% Awesome Cover Letter
% XeLaTeX Template
% Version 1.3 (30/3/2020)
%
% This template has been downloaded from:
% http://www.LaTeXTemplates.com
%
% Original authors:
% Claud D. Park (posquit0.bj@gmail.com)
% Lars Richter (mail@ayeks.de)
% With modifications by:
% Vel (vel@latextemplates.com)
%
% License:
% CC BY-NC-SA 3.0 (http://creativecommons.org/licenses/by-nc-sa/3.0/)
%
% Important note:
% This template must be compiled with XeLaTeX, the below lines will ensure this
%!TEX TS-program = xelatex
%!TEX encoding = UTF-8 Unicode
%
%%%%%%%%%%%%%%%%%%%%%%%%%%%%%%%%%%%%%%%%%

%----------------------------------------------------------------------------------------
%	PACKAGES AND OTHER DOCUMENT CONFIGURATIONS
%----------------------------------------------------------------------------------------

\documentclass[11pt, a4paper]{awesome-cv} % A4 paper size by default, use 'letterpaper' for US letter

\geometry{left=2cm, top=1.5cm, right=2cm, bottom=2cm, footskip=.5cm} % Configure page margins with geometry
 
\fontdir[fonts/] % Specify the location of the included fonts

% Color for highlights
\colorlet{awesome}{awesome-red} % Default colors include: awesome-emerald, awesome-skyblue, awesome-red, awesome-pink, awesome-orange, awesome-nephritis, awesome-concrete, awesome-darknight
%\definecolor{awesome}{HTML}{CA63A8} % Uncomment if you would like to specify your own color

% Colors for text - uncomment and modify
%\definecolor{darktext}{HTML}{414141}
%\definecolor{text}{HTML}{414141}
%\definecolor{graytext}{HTML}{414141}
%\definecolor{lighttext}{HTML}{414141}

\renewcommand{\acvHeaderSocialSep}{\quad\textbar\quad} % If you would like to change the social information separator from a pipe (|) to something else

%----------------------------------------------------------------------------------------
%	PERSONAL INFORMATION
%	Comment any of the lines below if they are not required
%----------------------------------------------------------------------------------------

\name{Marcellus E.}{Parker}
\address{101 Polk St., San Francisco, Ca 94102}
\mobile{(720)339-3093}

\email{mparker711@gmail.com}
\linkedin{linkedin.com/in/marcellus-parker}
%\homepage{www.posquit0.com}
%\github{posquit0}
%\linkedin{posquit0}
%\skype{skypeid}
%\stackoverflow{SOid}{SOname}
%\twitter{@twit}

%\position{Beamline Engineer - SSRL Materials Science Division} % Your expertise/fields
%\quote{``Make the change that you want to see in the world."} % A quote or statement

%----------------------------------------------------------------------------------------
%	RECIPIENT/POSITION/LETTER INFORMATION
%	All of the below lines must be filled out
%----------------------------------------------------------------------------------------

\recipient{SLAC National Laboratory}{2575 Sand Hill Road\\Menlo Park, CA 94025} % The company being applied to

\letterdate{\today} % The date on the letter, default is the date of compilation

\lettertitle{Job Application for: Beamline Engineer - SSRL Materials Science Division} % The title of the letter

\letteropening{Dear Hiring Manager,} % How the letter is opened

\letterclosing{Sincerely,} % How the letter is closed

%\letterenclosure[Attached]{Curriculum Vitae} % Any enclosures with the letter

\makecvfooter{\today}{Marcellus E. Parker~~~·~~~Cover Letter}{} % Specify the letter footer with 3 arguments: (<left>, <center>, <right>), leave any of these blank if they are not needed
  
%----------------------------------------------------------------------------------------

\begin{document}

\makecvheader % Print the header

\makelettertitle % Print the title

%----------------------------------------------------------------------------------------
%	LETTER CONTENT
%----------------------------------------------------------------------------------------

\begin{cvletter}

%------------------------------------------------




\lettersection{ About Me}
{\setlength{\parindent}{1.5em}
I am excited to apply for the Beamline Engineer position within SSRL’s Materials Science Division. I currently work as a Staff Engineer at SLAC, where I focus on restoring, optimizing, and supporting accelerator and diagnostic systems. This role strongly aligns with my interest in hands-on engineering, experimental systems, and user support—areas where I have built a solid foundation and continue to grow.

At SLAC, I have led the recommissioning of critical beamline systems, including the FACET positron beamline and a dormant bunch length diagnostic. These efforts required full ownership across vacuum systems, electrical infrastructure, cooling systems, and controls, with a focus on bringing complex, inactive systems back into reliable operation. In parallel, I have developed MATLAB-based tools for data acquisition and visualization to better understand and optimize beam performance, supporting both operations and experimental work.

Prior to SLAC, I worked at Fermilab as part of the US HL-LHC Accelerator Upgrade Project, where I focused on superconducting magnet R\&D, manufacturing, and electrical quality assurance. I worked directly on the fabrication, testing, and validation of Nb$_3$Sn superconducting magnet systems designed to operate under extreme thermal conditions, from cryogenic environments to high-temperature processing. I led electrical QA efforts, supported failure analysis, and collaborated closely with technicians and scientists to improve fabrication reliability and performance. This experience gave me a strong understanding of how complex systems are built, how they fail, and how to design and test them for robustness.
}

\lettersection{ Why I Fit}
{\setlength{\parindent}{1.5em}
Across both roles, I have worked in highly collaborative environments, interfacing regularly with physicists, engineers, technicians, and operations teams. I am comfortable taking ownership of systems, troubleshooting issues in real time, and supporting experimental needs in fast-paced settings. I am particularly drawn to this role at SSRL because of the opportunity to work more directly with beamline instrumentation and user-facing experimental systems while continuing to improve system reliability and performance.

While my background has been in accelerator beamlines, I am eager to apply that systems-level engineering experience to hard X-ray beamlines and experimental stations. I am confident that my combination of hands-on engineering, R\&D experience, and operational support would allow me to contribute effectively to SSRL from day one.
}

I would welcome the opportunity to discuss how my background aligns with the needs of your team.



%------------------------------------------------

%\lettersection{Why Initech?}

%Suspendisse commodo, massa eu congue tincidunt, elit mauris pellentesque orci, cursus tempor odio nisl euismod augue. Aliquam adipiscing nibh ut odio sodales et pulvinar tortor laoreet. Mauris a accumsan ligula. Class aptent taciti sociosqu ad litora torquent per conubia nostra, per inceptos himenaeos. Suspendisse vulputate sem vehicula ipsum varius nec tempus dui dapibus. Phasellus et est urna, ut auctor erat. Sed tincidunt odio id odio aliquam mattis. Donec sapien nulla, feugiat eget adipiscing sit amet, lacinia ut dolor. Phasellus tincidunt, leo a fringilla consectetur, felis diam aliquam urna, vitae aliquet lectus orci nec velit. Vivamus dapibus varius blandit.

%------------------------------------------------


%------------------------------------------------

\end{cvletter}

%----------------------------------------------------------------------------------------

\makeletterclosing % Print the signature and enclosures

\end{document}